\documentclass[11pt]{article}
\usepackage[margin=0.8in]{geometry}
\usepackage{amsmath,amssymb,booktabs,xcolor,framed}
\usepackage[colorlinks,linkcolor=blue!45!black]{hyperref}


\newcommand{\M}{\mathcal M}
\definecolor{bad}{RGB}{150,25,20}
\definecolor{good}{RGB}{20,100,55}
\definecolor{hd}{RGB}{40,40,60}

\newcommand{\LOC}[2]{\medskip\noindent\textcolor{hd}{\rule{\textwidth}{0.6pt}}\par
  \noindent\textbf{\large p.#1 \quad line #2}\par\smallskip}
\newcommand{\WRONG}[1]{\noindent\textcolor{bad}{\textbf{Currently reads:}}\par
  \noindent\begingroup\itshape #1\endgroup\par\smallskip}
\newcommand{\RIGHT}[1]{\noindent\textcolor{good}{\textbf{Replace with:}}\par
  \noindent\begingroup\itshape #1\endgroup\par\smallskip}
\newcommand{\WHY}[1]{\noindent\textbf{Why:}\ #1\par}

\setlength{\parindent}{0pt}
\setlength{\parskip}{3pt}

\title{Errata: line-by-line corrections\\\large \texttt{draft\_close\_to\_final\_aug28}}
\author{}
\date{28 August 2026}

\begin{document}
\maketitle

Every item below is: where it is, what it says now, what it should say. Line numbers are the
margin numbers of the reviewed PDF. Sorted by page.

Items marked \textbf{[DELETE]} require removing text, not replacing it.

% =====================================================================
\section*{Main paper}

\LOC{4}{155--157}
\WRONG{Therefore, $A<\cos\jmath$ is a sufficient condition for the existence of an interior
flow-versus-do-nothing threshold.}
\RIGHT{An interior threshold exists for every $A\in(0,1)$ and $\jmath\in(0,\pi/2)$: since
$\Delta(\sin\jmath)<0$ and $\Delta(1)>0$, a root always lies in $(\sin\jmath,1)$. The condition
$A<\cos\jmath$ characterises \emph{uniqueness} of the transition: it holds if and only if
$\Delta$ has a single sign change on $(0,1)$.}
\WHY{The stronger statement is proved in \texttt{threshold.tex} and verified on 120 cells with no
exceptions. Theorem~2's statement should be updated to match.}

\LOC{6}{207--209}
\WRONG{\dots\ 0.338 (predicted flow), 0.322 (learned flow), 0.264 (do nothing) \dots}
\RIGHT{\dots\ 0.338 (predicted flow), 0.3184 (learned flow), 0.264 (do nothing) \dots}
\WHY{Line 217 of the same section, Table~C4 and Table~B3 all give $\M_1=0.3184$. The value
$0.322$ appears only here and in Table~C1. Make one choice and use it everywhere.}

\LOC{6}{209}
\WRONG{\dots\ and 0.137 (perfect sampler).}
\RIGHT{\dots\ and 0.137 (perfect sampler, an independent draw sharing the posterior mean; a draw
conditioned on $\bm c$ alone scores 0.0500).}
\WHY{Two distinct reference quantities exist, $A^2=0.1369$ and $(As)^2=0.0493$; the measured
true-posterior row is $0.0500$. As written the reader cannot tell which is meant.}

\LOC{7}{227--228}
\WRONG{\dots\ the gap is 0.0985 at width 256 and 0.1040 at width 2048.}
\RIGHT{\dots\ the gap remains positive in all twelve learning-rate by width cells, from 0.0419 to
0.1158, with the same sign in every cell at a minimum of 16 standard errors (Table~B2).}
\WHY{$0.0985/0.1040$ come from the earlier single-learning-rate sweep (Table~B3), which is a
different experiment from the three-rate grid in Table~B2. The two tables are not reconcilable;
keep B2 and delete B3.}

\LOC{7}{242--243}
\WRONG{Increasing integration steps from one to sixty-four lowers terminal cosine from 0.350 to
0.319, while mean pairwise cosine decreases from 0.5299 to 0.1829 \dots}
\RIGHT{Increasing integration steps from one to sixty-four lowers terminal cosine from 0.3541 to
0.3172, while mean pairwise cosine decreases from 0.5299 to 0.1829 \dots}
\WHY{The pairwise-cosine values are from the current run; the cosine values are from an earlier
one. Table~B6 gives 0.3541 and 0.3172.}

\LOC{7}{252--253}
\WRONG{\dots\ the flow increases target alignment from 0.50 to 0.69, compared with 0.79 for the
corresponding class prototype.}
\RIGHT{\dots\ the flow increases target alignment from 0.5517 to 0.7586, compared with 0.7930 for
the corresponding class prototype and a random-pair floor of 0.4815.}
\WHY{Two independent runs, the sweep and the probe rerun, both give these values. The $0.50/0.69$
pair is from an early pass. Adding the floor is what makes the gain interpretable.}

\LOC{7}{264}
\WRONG{At both completed resolutions, the source-conditioned model $\M_2$ achieves slightly
higher PSNR \dots}
\RIGHT{At all three resolutions, the source-conditioned model $\M_2$ achieves equal or slightly
higher PSNR \dots}
\WHY{$96\times96$ is complete. There $\M_1$ and $\M_2$ score 23.1984 and 23.1984 --- identical to
four decimals --- so ``slightly higher'' is wrong at that resolution. Equal score with an
18\% worse energy distance is the strongest form of the claim; it should be stated, not dropped.}

\LOC{8}{303--305}
\WRONG{Finally, the $1/t$ source-recovery singularity is specific to the interpolation analyzed
here; whether the same mechanism persists under alternative paths, stochastic dynamics, and
high-dimensional image-space flow matching remains open.}
\RIGHT{The $1/t$ source-recovery singularity is derived for the linear interpolant; it carries
over to the geodesic interpolant used in the experiments, where the recovery sensitivity is
$\sin((1-\alpha t)\theta)/\sin((1-\alpha)t\theta)$ and diverges at the same rate, differing only
by the bounded factor $\sin\theta/\theta$ (Appendix~A.6). A discrete masked-diffusion analogue,
reported in Appendix~B, does \emph{not} reproduce the source-conditioning penalty: there,
conditioning on the source while starting from it helps by $+0.0107$ on average. The mechanism is
therefore specific to the continuous deterministic setting.}
\WHY{Both facts are established and currently absent. Leaving the question ``open'' when it was
tested and the sign reversed is the weaker and riskier option.}

% =====================================================================
\section*{Appendix A}

\LOC{12}{427--435}
\WRONG{Three \texttt{[TODO---Yassine]} blocks requesting the marginal-similarity parameter, the
regression, the asymptotic-slope derivation, and the first-order perturbation derivation.}
\RIGHT{[DELETE the three blocks.] Insert the content of \texttt{ratio.tex}: the sweep parameter
is $\delta$ in $\Sigma_1(\delta)=\Sigma_0+\delta(\Sigma_1-\Sigma_0)$ with tested values
$0.05, 0.1, 0.2, 0.4, 0.7, 1.0$ and corresponding relative marginal separations
$0.043, 0.087, 0.173, 0.347, 0.607, 0.867$; the regression is ordinary least squares of
$\log_{10}$ ratio on $\log_{10}\delta$ over those six points, slope $-0.756$, rms residual
$0.0166$; the asymptotic slope $-1$ follows because the span response carries a factor that is
$O(\delta)$ while the skew response is $O(1)$.}
\WHY{All answered. See also the next item, which corrects the explanation currently given for the
shortfall.}

\LOC{12}{426}
\WRONG{Across these marginal configurations, the fitted log--log slope is $-0.756$, rather than
the reported asymptotic prediction $-1$.}
\RIGHT{Across these marginal configurations the fitted log--log slope is $-0.756$. The ratio
slope is identically the skew slope minus the span slope. Over this window the span exponent is
$0.573$ and the skew exponent $-0.159$; a decade lower they are $0.933$ and $-0.019$, giving a
ratio slope of $-0.952$. The shortfall from $-1$ is therefore the span response being sublinear
over a pre-asymptotic window, not the decline of the skew constant, which pushes the slope
\emph{toward} $-1$.}
\WHY{If the paper attributes the shortfall to the skew constant, the arithmetic contradicts it:
that term contributes $-0.16$ in the opposite direction.}

\LOC{14}{498--502}
\WRONG{\texttt{[TODO---Yassine]} asking for a proof that $\Delta$ has no root in $(0,1)$ when
$A\ge\cos\jmath$.}
\RIGHT{[DELETE the block.] Insert \texttt{threshold.tex}. There is no such proof because the
statement is false: $\Delta(\sin\jmath)<0$ for every $A<1$, so a root always exists in
$(\sin\jmath,1)$. What $A<\cos\jmath$ characterises is uniqueness.}
\WHY{See p.4 l.155 above. This also fills three cells of Table~B1.}

\LOC{15}{521--545}
\WRONG{Section A.6 derives the recovery relation and the $1/t$ sensitivity for
$z_t=(1-t)z_0+tz_1$ only.}
\RIGHT{[ADD after A.6.8] On the geodesic interpolant used in all experiments, $z_t$ and the
auxiliary state both lie in $\mathrm{span}\{z_0,z_1\}$; the $2\times2$ recovery determinant
collapses to $-\sin((1-\alpha)t\theta)/\sin\theta$ and the sensitivity is
$\sin((1-\alpha t)\theta)/\sin((1-\alpha)t\theta)$, which diverges as $1/t$ exactly as in the
flat case, differing by the bounded factor $\sin\theta/\theta\in(0,1]$; the remainder is second
order in $t$ when $\alpha=0$.}
\WHY{Every experiment uses the geodesic while every proposition in \S6 assumes the linear
interpolant. Verified against finite differences in dimension 64 to $2.7\times10^{-10}$.}

% =====================================================================
\section*{Appendix B}

\LOC{16}{571--573}
\WRONG{\dots\ the measured threshold agrees with the analytic prediction to a maximum absolute
error of $3.3\times10^{-2}$.}
\RIGHT{\dots\ the measured threshold agrees with the analytic prediction to a maximum absolute
error of $3.3\times10^{-2}$. The discrepancy is one-sided: the measured value exceeds the
analytic one in 22 of the 27 cells, with mean signed error $+0.0099$ (sign test
$p=1.5\times10^{-3}$), consistent with the finite fitting accuracy of the oracle used to
integrate the population field.}
\WHY{A reader who plots the comparison sees the bias immediately. Reporting only the maximum
absolute error looks like it was not noticed.}

\LOC{16}{Table B1}
\WRONG{Entries ``--'' at $(A,\jmath)=(0.70,0.80)$, $(0.90,0.50)$ and $(0.90,0.80)$.}
\RIGHT{$0.8347$, $0.5869$ and $0.7708$ respectively, with measured values $0.8286$, $0.5600$ and
$0.7586$.}
\WHY{These cells have two roots, not none; the sweep already measured the upper one, and it
agrees with the analytic value to $\le0.027$. The comparison becomes 30 cells rather than 27,
and the sentence at l.571 should be updated accordingly.}

\LOC{18}{Table B3}
\WRONG{Width / $\M_1$ / $\M_2$ / Gap table with 0.3138, 0.2153, 0.0985 and 0.2904, 0.1864,
0.1040.}
\RIGHT{[DELETE the table.] Fold its training-loss and distributional-error trend (l.598--603)
into the discussion of Table~B2.}
\WHY{Table~B3 is the earlier single-learning-rate sweep and Table~B2 is the three-rate grid.
Keeping both invites the question of why the same gap has two values.}

\LOC{19}{610--616}
\WRONG{Two \texttt{[TODO---Yassine]} blocks on the meaning of ``$4\times$ training'' and on
whether the capacity sweep was repeated at a second learning rate.}
\RIGHT{[DELETE both blocks.] $4\times$ means optimisation steps, $3{,}000$ against $12{,}000$,
with batch size, optimizer, learning rate, parameterisation, gradient clipping, checkpoint
selection, evaluation protocol and seeds all unchanged. The sweep has been repeated at three
learning rates, $10^{-4}$, $3\times10^{-4}$ and $10^{-3}$; see Table~B2.}

\LOC{20}{658--664}
\WRONG{The two interventions are compared without stating that they were originally sampled with
different integrators.}
\RIGHT{[ADD] At $\alpha=0$ and $\sigma_{\mathrm{aux}}=0$ both auxiliary families use the same
trained model (terminal losses $4.9116\times10^{-4}$ and $4.9046\times10^{-4}$). Sampled with
different integrators they reported distributional errors of $0.0782$ and $0.0367$; re-run under
a single integrator the anchor agrees and the two families trace one curve, with a
cross-to-within ratio of $1.40$ against $3.44$ under mismatched integrators.}
\WHY{Without this, a reader comparing the two families is comparing two integrators.}

\LOC{21}{Evaluation protocol, B.5}
\WRONG{No statement of evaluation noise.}
\RIGHT{[ADD] The distributional report fixes its evaluation seed; the fidelity report does not.
Scoring one identical trained network twice therefore yields cosines differing by about $0.002$,
which sets a floor below which differences in any reported cosine or pairwise cosine are not
meaningful. All gaps reported in this work are between $0.04$ and $0.12$.}
\WHY{This resolves the several places where the same model appears with slightly different
cosines, without editing any of them.}

\LOC{22}{719--721}
\WRONG{\dots\ increasing the guidance weight changes terminal target cosine by $+0.157$ under
noise initialization, but by $-0.505$ under informative initialization.}
\RIGHT{\dots\ increasing the guidance weight changes terminal target cosine by
$+0.0770_{\pm0.0004}$ under noise initialization, but by $-0.4905_{\pm0.0039}$ under informative
initialization.}
\WHY{The main text at l.234 and Table~C6 already use the corrected values; only this paragraph
still carries the old pair.}

\LOC{22}{722}
\WRONG{\dots\ present in the initialization. \texttt{:contentReference[oaicite:0]index=0}}
\RIGHT{\dots\ present in the initialization.}
\WHY{[DELETE] Paste artefact.}

\LOC{22}{735--739}
\WRONG{\texttt{[TODO---Yassine]} asking for the guidance formulation, weights, $\lambda$, the
terminal-cosine values and the unguided reference.}
\RIGHT{[DELETE the block.] The formulation is
$\bm v_w=\bm v_{\mathrm{uncond}}+w(\bm v_{\mathrm{cond}}-\bm v_{\mathrm{uncond}})$ with tangent
projection, $w\in\{1,2,3\}$, $10\%$ condition dropout, and $w=1$ as the unguided reference;
$\lambda$ interpolates the start via $\mathrm{slerp}(\bm u,\bm z_0,\lambda)$ with tested values
$0,0.25,0.5,0.75,1$; the full table is Table~C6.}

\LOC{22}{746--748}
\WRONG{The blurred-source embedding has target alignment 0.50, while applying the flow increases
the alignment to 0.69. For comparison, the corresponding class prototype attains 0.79.}
\RIGHT{The blurred-source embedding has target alignment 0.5517, while applying the flow
increases the alignment to 0.7586. For comparison, the corresponding class prototype attains
0.7930 and a random pair of embeddings scores 0.4815, so the flow recovers 89\% of the available
range.}

\LOC{22}{749--751}
\WRONG{The measured factorization discrepancy is 0.62, indicating substantial departure from the
decomposition assumed in the synthetic analysis.}
\RIGHT{The estimated parameters are $\hat A=0.7923$, $\hat\beta=0.6010$ and
$\hat\jmath=0.2817$, giving a predicted threshold of $0.5291$; $\hat\beta$ lies above it, so the
prediction that the flow beats the source was made in advance and holds.}
\WHY{No quantity named a factorization discrepancy is computed anywhere. The stage was re-run to
check and no output column equals $0.62$; the nearest value is $\hat\beta=0.6010$, which already
appears under its own name. The sentence cannot be sourced.}

\LOC{23}{Table B7}
\WRONG{Blurred source 0.50 \quad Flow output 0.69 \quad Class prototype 0.79}
\RIGHT{Random pair 0.4815 \quad Blurred source 0.5517 \quad Flow output 0.7586 \quad
Source-conditioned 0.7507 \quad Class prototype 0.7930}
\WHY{Adding the floor and the source-conditioned column shows that the model ordering also
transfers to the real embedding space, which is a result the table currently omits.}

\LOC{23}{754--757}
\WRONG{\texttt{[TODO---Yassine]} asking for the OpenCLIP architecture, split, preprocessing, blur,
normalisation, metric, prototype and the $0.62$.}
\RIGHT{[DELETE the block.] OpenCLIP ViT-B-32, checkpoint \texttt{laion2b\_s34b\_b79k}; CIFAR-10
train split, 6{,}000 images at a fixed seed; the checkpoint's own preprocessing; $L_2$-normalised
embeddings; Gaussian blur with kernel 9 and $\sigma\sim\mathrm{Unif}[0.8,3.0]$; condition is the
normalised class-mean of clean embeddings; $80/20$ split; width 512, 3{,}000 steps at
$\eta=10^{-4}$, 32 Heun steps; target alignment is the mean cosine to the clean embedding of the
same image on the held-out 20\%. On the $0.62$, see p.22 l.749 above.}

\LOC{23}{765--768}
\WRONG{\texttt{[TODO---Yassine]} asking for the completed super-resolution experiment.}
\RIGHT{[DELETE the block.] The experiment is complete at all three resolutions; the content is in
Table~C7 and Table~C8, which move here when Appendix~C is merged.}

% =====================================================================
\section*{Appendix C --- to be deleted}

\LOC{24}{769--959}
\WRONG{Appendix~C repeats Appendix~B: threshold, reference ladder, guidance, capacity, CLIP and
super-resolution each appear twice, with different numbers.}
\RIGHT{[DELETE Appendix~C.] Move Tables C1--C9 and Figure~5 into Appendix~B, replacing the
corresponding B-tables where C is the newer version. Specifically: C6 (guidance) supersedes the
B.6 text; C7 and C8 (super-resolution) fill B.8; C1 (reference ladder) belongs in B.4.}
\WHY{This is the single largest source of internal contradiction, and iThenticate compares the
document against itself, so duplicated appendices raise the similarity score directly.}

\LOC{24}{790--791}
\WRONG{Using one integration step gives terminal cosine similarity 0.350, whereas using sixty-four
steps gives 0.319 \dots}
\RIGHT{Using one integration step gives terminal cosine similarity 0.3541, whereas using
sixty-four steps gives 0.3172 \dots}
\WHY{Table~B6 in the same document gives the corrected values.}

\LOC{27}{881--882}
\WRONG{\dots\ terminal fidelity increases from 0.2195 at zero perturbation to 0.3122 at the
strongest tested perturbation, compared with 0.3195 for $\M_1$.}
\RIGHT{\dots\ terminal fidelity increases from 0.2149 at zero perturbation to 0.3105 at the
strongest tested perturbation, compared with 0.3180 for $\M_1$.}
\WHY{Table~B5 in the same document gives 0.2149 and 0.3105.}

\LOC{27}{885--893}
\WRONG{The CLIP paragraph does not report the label probe.}
\RIGHT{[ADD] A linear probe trained on the component of the source embedding orthogonal to the
condition reaches 0.058 accuracy on the held-out split, against a chance level of 0.10, and
therefore does not extract label information from that component.}
\WHY{An earlier in-sample version of this probe was withdrawn as instrumentation; the held-out
version restores the result and should be reported as such.}

\LOC{28}{894--895}
\WRONG{\dots\ the flow increases target alignment from 0.50 to 0.69, compared with 0.79 \dots}
\RIGHT{\dots\ the flow increases target alignment from 0.5517 to 0.7586, compared with 0.7930
\dots}

\LOC{28}{922}
\WRONG{At both completed resolutions \dots}
\RIGHT{At all three resolutions \dots}

\LOC{28}{Table C8}
\WRONG{Two resolution blocks, $32\times32$ and $48\times48$.}
\RIGHT{[ADD a third block] Resolution $96\times96$: identity PSNR 21.0806, energy 17.1672,
$W_1^{\mathrm{hf}}$ 998.74; $\M_1$ 23.1984, 1.8884, 623.56; $\M_2$ 23.1984, 2.2222, 659.37; null
floor 0.0010 and 17.63.}

\LOC{28}{927--929}
\WRONG{\texttt{[TODO---Yassine]} asking for the sample count, seeds, preprocessing and the
$96\times96$ results.}
\RIGHT{[DELETE the block.] All arms are evaluated on the full STL-10 test split of 1{,}024
images with three training seeds; the loader is the same function used by the dimensionality
measurement, so the preprocessing and split are identical by construction; the $96\times96$
results are above.}

\LOC{29}{942--944}
\WRONG{\texttt{[TODO---Yassine]} asking how the seven strengths are parameterised and what
``retaining'' versus ``transporting'' means.}
\RIGHT{[DELETE the block.] The strength is the image-to-image pipeline's own \texttt{strength}
argument, taking the seven values 0.1, 0.2, 0.3, 0.4, 0.5, 0.6 and 0.8: the degraded image is
encoded to a latent, noised to a fraction $s$ along the schedule and denoised from there, so
$\lceil 30s\rceil$ of 30 steps are executed. Retaining means returning the degraded image
unchanged, and is the baseline with gain zero by definition; transporting means returning the
pipeline output at strength $s$. Both are scored by CLIP cosine to the clean original.}

% =====================================================================
\section*{Figures}

\LOC{6}{Figure 1}
\WRONG{Three panels at one-third column width; axis labels and legends unreadable at print size.}
\RIGHT{Use the three panels as separate figures, or give the composite full page width. Panel~(a)
must label the green points ``population field (measured)'', not ``flow'': they are the fitted
oracle for the population field, which is why they sit near 0.335 rather than $\M_1$'s 0.3184.}

\LOC{17}{Figure 2}
\WRONG{Plots the $j=0$ points, which Table~B1's own caption calls formal only, and the three
$A\ge\cos\jmath$ cells, which is why isolated points appear at $A=0.9$ with no curve through
them.}
\RIGHT{Either restrict to the 27 cells with a predicted threshold, or adopt the corrected
Theorem~2 and plot all 30 with their analytic upper roots, which gives every point a curve.}

\LOC{18}{Figure 3}
\WRONG{Legend entries are malformed: ``source-conditioned $\eta$:0.000 1''.}
\RIGHT{Regenerate with the legend labels fixed.}

\LOC{---}{No qualitative figure anywhere}
\WRONG{The paper contains no images despite two image experiments.}
\RIGHT{[ADD] Two panels are available. The super-resolution panel (degraded input, regression,
$\M_1$, $\M_2$, target) shows $\M_2$ visibly smoother than $\M_1$ while scoring higher, and the
regression arm sharpest of all while scoring worst. The image-to-image panel shows, per
degradation level, the clean target, the retained source, and the output at three strengths.}
\WHY{The central claim of the paper --- higher score, less restored detail --- is a visual
property. A reader currently has to take it on trust.}

\end{document}
